\section{Conclusion}
\label{sec:conclusion}

For a bound Newtonian two-body orbit, the geocentric position of a body
is exactly a sum of two Bessel--Fourier epicycle families whose radii,
planes and rates are determined by the orbital elements. We gave the explicit
three-dimensional expansion with truncation bounds and verified them: $N\approx6$ to $10$ harmonics, i.e.\ $2N=12$ to $20$ rotating circles
plus an offset per orbit, reach $10^{-6}$ to $10^{-9}$~AU for
Mars and Ceres, hundreds are needed at $e=0.89$, and a circular two-term approximation
of Mars (not Ptolemy's actual model) reproduces the retrograde loops but is off by up to
$33^\circ$. Frequency analysis recovers the analytic epicycles from
Keplerian data to $\sim10^{-10}$ in frequency and, for a perturbed Ceres
trajectory, represents the interior of a $200$-yr window with errors below
$10^{-4}$~AU at $K=200$ while remaining edge-limited over the whole window,
where it is no better than a Kepler ellipse fitted to the same heliocentric data with matched weights. The right
statement of the hypothesis is therefore bounded and finite-window:
exact for Kepler motion, almost-periodic for the geocentric view,
exact on analytic invariant tori, an empirical finite-window approximation for
other perturbed motion, and absent for unbound orbits. The accompanying interactive viewer makes the
convergence visible. Natural next steps are validation against a modern
ephemeris, a treatment of edge effects by better windows or extrapolation, and
extension to more bodies.
